\chapter{Power Markets II: Trading}\label{ch:m3:power-markets-trading}

Forty-five minutes before four o'clock, a wind farm's forecast for the next hour falls
by a third. It sold 100 megawatt-hours for that hour in yesterday's auction; now it
expects to produce 70. It can buy the missing 30 on the continuous \omterm{def:m3:power-markets-trading:intraday}{intraday market}, at
whatever price the order book shows, or deliver short and let the grid operator charge
it the \omterm{def:m3:power-markets-trading:imbalance}{imbalance price}, known only weeks later. Most of the time the difference is a
few euros. On a tight evening it can be hundreds. Power trading after the day-ahead
auction is the business of managing that difference: forecasts that move, a market that
closes shortly before delivery, and a settlement that punishes whoever is wrong. This
chapter covers the \omterm{def:m3:power-markets-trading:intraday}{intraday market}, imbalance settlement, the instruments that hedge
congestion, the economics of renewable output and its long-term contracts, and batteries.

\section{Continuous intraday markets}

\begin{definition}[Intraday market, gate closure]\label{def:m3:power-markets-trading:intraday}
An \emph{intraday market}\index{intraday market} trades power for delivery later the same
day or the next day, after the day-ahead auction, either continuously in an order book with
price--time priority (One Quant Book 1, chapter 19) or in intraday auctions. Its \emph{gate
closure}\index{gate closure} for an MTU is the last moment at which that MTU can be traded;
after it, deviations are settled as imbalances.
\end{definition}

Intraday trading exists because forecasts improve as delivery approaches: a wind forecast a
day ahead may be off by a fifth of the output, an hour ahead by much less. Each producer and
supplier corrects its position as its forecast moves, and prices move with the aggregate
forecast error of the zone.

\begin{dated}{2026-09}{European intraday coupling}\label{dat:m3:power-markets-trading:sidc}
The Single Intraday Coupling (SIDC), launched on 12--13 June 2018 in fifteen countries,
lets orders in one zone match orders in another while cross-border capacity is free. Cross-zonal
trading closed 60 minutes before the start of each MTU; a 30-minute cross-zonal \omterm{def:m3:power-markets-trading:intraday}{gate closure}
has since gone live on some borders. Trading within a zone can continue closer to delivery,
depending on the country.
\end{dated}

\begin{omfigure}
\centering
\begin{tikzpicture}[font=\small]
\draw[thick,-{Stealth[length=2.5mm]}] (0,0) -- (13.2,0);
\foreach \x/\t/\l in {0.4/{months ahead}/{forwards and\\futures},3.2/{D$-1$, 12:00}/{day-ahead\\auction},6.4/{D$-1$ afternoon\\to delivery}/{continuous\\intraday},9.4/{gate closure}/{last trade},12.0/{weeks later}/{imbalance\\settlement}}{
  \draw[thick] (\x,0.12) -- (\x,-0.12);
  \node[anchor=south,align=center] at (\x,0.2) {\t};
  \node[anchor=north,align=center] at (\x,-0.2) {\l};}
\draw[omThm,thick] (10.7,0.3) -- (10.7,-0.3);
\node[omThm,anchor=north,font=\footnotesize] at (10.7,-1.0) {delivery};
\end{tikzpicture}

\omcaption{The markets for one MTU, from the forward market to settlement. Every trade before
\omterm{def:m3:power-markets-trading:intraday}{gate closure} changes the schedule; every deviation after it is an imbalance. Schematic.}\label{fig:m3:power-markets-trading:timeline}
\end{omfigure}

\section{Imbalance settlement and balancing responsibility}

\begin{definition}[Balancing responsible party]\label{def:m3:power-markets-trading:brp}
A \emph{balancing responsible party}\index{balancing responsible party} (BRP) is a market
participant, or its representative, that is financially responsible for the difference between
the energy it scheduled and the energy it actually injected or withdrew, for a portfolio of
connection points, in each imbalance settlement period.
\end{definition}

\begin{definition}[Imbalance volume, imbalance price]\label{def:m3:power-markets-trading:imbalance}
The \emph{imbalance volume}\index{imbalance volume} of a BRP in a settlement period is its
metered net injection minus its scheduled net injection: positive when it is long (delivered
more than it sold), negative when short. The \emph{imbalance price}\index{imbalance price} is
the price at which the grid operator settles it, derived from the cost of the balancing energy
the operator activated in that period.
\end{definition}

Under a single price, a long BRP is paid the \omterm{def:m3:power-markets-trading:imbalance}{imbalance price} for its surplus and a short one
pays it for its deficit. When the system as a whole is short, the operator activates upward
reserves and the \omterm{def:m3:power-markets-trading:imbalance}{imbalance price} is high: being short then is expensive, and being long is
rewarded. The \omterm{def:m3:power-markets-trading:imbalance}{imbalance price} is therefore a penalty only for those who move in the same
direction as the system, and a BRP whose error has the opposite sign earns from it. Most
European systems settle imbalances per 15 minutes; Germany computes one uniform price per
quarter-hour for all its balancing groups.

\begin{example}[A wind farm's afternoon]\label{ex:m3:power-markets-trading:wind}
A wind farm sells 80, 90, 100 and 100~MWh day-ahead for four hours at 62, 58, 55 and
\qty{57}{EUR/MWh}. Before the last hour its forecast falls by 30~MWh and it buys 30 back at
\qty{170}{EUR/MWh} intraday. It meters 82, 88, 101 and 55~MWh; the \omterm{def:m3:power-markets-trading:imbalance}{imbalance prices} are 60,
75, 50 and \qty{240}{EUR/MWh}. Its trading revenue is EUR~16{,}280 (the day-ahead sales less
the buy-back), and its imbalances ($+2$, $-2$, $+1$ and $-15$~MWh) cost it EUR~3{,}580, nearly
all in the last hour, when it was short while the system was short too: EUR~12{,}700 in all.
\end{example}

\section{Transmission rights and virtual bids}

In nodal markets a generator is paid its node's price and a supplier pays its load's, so
congestion between them is a risk the market lets them hedge.

\begin{definition}[Financial transmission right]\label{def:m3:power-markets-trading:ftr}
A \emph{financial transmission right}\index{financial transmission right} (FTR) from node $A$
to node $B$ entitles its holder, for each hour of its term, to a quantity times the day-ahead
congestion price at $B$ minus that at $A$: a payment when the path is congested in the
direction of the right, and, for an obligation, a charge when it is congested the other way.
It is sold by the grid operator in auctions and funded by the \omterm{def:m3:power-markets-design:lmp}{congestion rents} it collects.
\end{definition}

A generator at $A$ selling to a load at $B$ at a fixed price owns the congestion risk between
them; an FTR from $A$ to $B$ turns it into a fixed auction price.

\begin{definition}[Virtual bid]\label{def:m3:power-markets-trading:virtual}
A \emph{virtual bid}\index{virtual bid} is a purely financial position in a two-settlement
market: an offer to sell (an increment offer) or a bid to buy (a decrement bid) at a node in
the \omterm{def:m3:power-markets-design:da}{day-ahead market}, liquidated automatically at the real-time price, with no physical
delivery.
\end{definition}

A trader who expects real-time prices at a node to exceed day-ahead prices buys day-ahead with
a decrement bid and is paid the difference; the market wants this, because it pulls the
day-ahead price towards the expected real-time price. The same trades can manipulate: a
virtual that loses money on its own but raises the value of the trader's FTRs on the same path
is the classic abuse, and PJM's rules withhold FTR payments in such cases.

\section{Renewable output, capture prices and power purchase agreements}

A wind farm or a solar park does not choose when to produce. What it earns on the spot
market is not the average price but the average weighted by its own output.

\begin{definition}[Capture price]\label{def:m3:power-markets-trading:capture}
The \emph{capture price}\index{capture price} of a generator over a period is the average of
the market price weighted by its output, $\sum_h p_h q_h / \sum_h q_h$; its capture rate is
the capture price divided by the time-weighted average (baseload) price.
\end{definition}

Solar parks in one zone all produce in the same hours, which are the hours they push the
price down (\cref{ch:m3:power-markets-design}): the more solar is built, the lower its capture
rate. In the German data, solar's \omterm{def:m3:power-markets-trading:capture}{capture price} was \qty{46.23}{EUR/MWh} in 2024 against a
baseload average of \qty{78.51}{} (a capture rate of 59\%), and \qty{46.07}{} against
\qty{89.32}{EUR/MWh} in 2025 (52\%); onshore wind captured 82\% and 87\%
(\cref{fig:m3:power-markets-trading:capture}).

\begin{omfigure}
\begin{tikzpicture}
\begin{axis}[width=11cm,height=4.8cm,
  ylabel={EUR/MWh},
  tick label style={font=\small},label style={font=\small},
  xmin=2024,xmax=2026,ymin=0,ymax=160,grid=major,grid style={gray!25},
  xtick={2024,2024.5,2025,2025.5,2026},xticklabels={Jan 2024,Jul 2024,Jan 2025,Jul 2025,Jan 2026},
  legend columns=2,legend style={font=\footnotesize,at={(0.5,-0.2)},anchor=north,draw=none,fill=none,/tikz/every even column/.append style={column sep=8pt}},
  table/col sep=comma]
\addplot[omDef,very thick,mark=*,mark size=1.2pt] table[x=t,y=base]{figdata/markets-3/06-power-markets-trading/capture.csv};
\addplot[omProp,very thick,mark=square*,mark size=1.2pt] table[x=t,y=solar]{figdata/markets-3/06-power-markets-trading/capture.csv};
\legend{baseload average,solar capture price}
\end{axis}
\end{tikzpicture}

\omcaption{Germany, monthly: the average day-ahead price and the price captured by solar
output, 2024--2025. The gap is widest in the sunniest months. Data: Bundesnetzagentur |
SMARD.de, CC BY 4.0.}\label{fig:m3:power-markets-trading:capture}
\end{omfigure}

\begin{definition}[Power purchase agreement, shape risk]\label{def:m3:power-markets-trading:ppa}
A \emph{power purchase agreement}\index{power purchase agreement} (PPA) is a long-term contract
under which a buyer, often a large consumer, buys a generator's output (or a financial
equivalent) at an agreed price for years. \emph{Shape risk}\index{shape risk} is the risk
that the profile of the delivered output differs in value from a baseload block of the same
volume: in a pay-as-produced PPA it falls on the buyer, in a baseload PPA on the seller, who
must buy in the hours it does not produce.
\end{definition}

A pay-as-produced solar PPA at \qty{60}{EUR/MWh} was a good hedge for a buyer in 2024 only if
it valued the output at its \omterm{def:m3:power-markets-trading:capture}{capture price}: at a \omterm{def:m3:power-markets-trading:capture}{capture price} of 46, the buyer paid about
EUR~14 a MWh above the market value of what it received.

\begin{dated}{2026-09}{How much is contracted and stored}\label{dat:m3:power-markets-trading:ppa}
Pexapark counted about 6.08~GW of renewable capacity contracted under PPAs in Europe in the first half of 2025, in 124 deals, 26\% less than a
year earlier. In Texas, a presentation to ERCOT's Technology and Security Committee in February 2026 put installed and operating battery capacity at
15{,}712~MW, over 6{,}000~MW of it added in the previous year.
\end{dated}

\section{Batteries and flexible assets}

\begin{definition}[Round-trip efficiency]\label{def:m3:power-markets-trading:rte}
The \emph{round-trip efficiency}\index{round-trip efficiency} of a storage asset is the energy
it returns divided by the energy it took to charge it; the rest is lost in conversion.
\end{definition}

A battery buys cheap hours and sells dear ones. With perfect foresight of day-ahead prices,
its best schedule each day is a small dynamic programme over its state of charge, a problem of
the kind One Quant Book 4, chapter 9, treats in general.

\begin{method}[Day-ahead arbitrage of a battery]\label{met:m3:power-markets-trading:battery}
Discretise the state of charge in steps of one hour at full power. For each hour, from the
first, keep for each reachable state (charge level, cycles used) the best cash so far; move by
charging (pay the price for $1/\sqrt{\eta}$ times the energy stored), discharging (receive the
price for $\sqrt{\eta}$ times the energy released) or idling; at the end of the day take the
best empty state.
\end{method}

A 100~MW, 200~MWh battery with a \omterm{def:m3:power-markets-trading:rte}{round-trip efficiency} of 88\% and one full cycle a day
(illustrative figures) would have earned, with perfect foresight of the German day-ahead
prices, EUR~65{,}968 per MW in 2024 and EUR~74{,}546 per MW in 2025; allowing two cycles a day
raises 2025 to EUR~88{,}688. Perfect foresight overstates what a real operator earns from the
auction alone, while \omterm{def:m3:power-markets-trading:intraday}{intraday markets} and reserves add revenues this calculation omits.

\begin{omfigure}
\begin{tikzpicture}
\begin{axis}[width=11cm,height=4.6cm,axis y line*=right,axis x line=none,
  ylabel={state of charge, MWh},ylabel near ticks,
  tick label style={font=\small},label style={font=\small},
  xmin=-0.5,xmax=23.5,ymin=0,ymax=440,ytick={0,100,200},
  table/col sep=comma]
\addplot[omDef,thick,const plot,fill=omDef!12] table[x=tend,y=soc]{figdata/markets-3/06-power-markets-trading/batteryday.csv}\closedcycle;\label{plot:m3:battsoc}
\end{axis}
\begin{axis}[width=11cm,height=4.6cm,axis y line*=left,
  xlabel={hour, German local time},ylabel={EUR/MWh},
  tick label style={font=\small},label style={font=\small},
  xmin=-0.5,xmax=23.5,ymin=-20,ymax=200,grid=major,grid style={gray!25},
  xtick={0,3,6,9,12,15,18,21},
  legend columns=2,legend style={font=\footnotesize,at={(0.5,-0.33)},anchor=north,draw=none,fill=none,/tikz/every even column/.append style={column sep=8pt}},
  table/col sep=comma]
\addplot[omThm,very thick,const plot mark mid] table[x=hour,y=price]{figdata/markets-3/06-power-markets-trading/batteryday.csv};\addlegendentry{day-ahead price}
\addlegendimage{/pgfplots/refstyle=plot:m3:battsoc}\addlegendentry{state of charge (end of hour)}
\end{axis}
\end{tikzpicture}

\omcaption{The battery's best day-ahead schedule on 18 June 2025: it charges in the two
cheapest hours after noon (prices just below zero) and discharges in the two dearest of the
evening, earning EUR~30{,}149. Illustrative battery; prices: Bundesnetzagentur | SMARD.de,
CC BY 4.0.}\label{fig:m3:power-markets-trading:battery}
\end{omfigure}

\begin{omfigure}
\begin{tikzpicture}
\begin{axis}[width=11cm,height=4.6cm,
  xlabel={day of the year},ylabel={thousand EUR per MW},
  tick label style={font=\small},label style={font=\small},
  xmin=1,xmax=366,ymin=0,ymax=80,grid=major,grid style={gray!25},
  legend columns=2,legend style={font=\footnotesize,at={(0.5,-0.3)},anchor=north,draw=none,fill=none,/tikz/every even column/.append style={column sep=8pt}},
  table/col sep=comma]
\addplot[omProp,very thick] table[x=doy,y=y2024]{figdata/markets-3/06-power-markets-trading/batterycum.csv};
\addplot[omDef,very thick] table[x=doy,y=y2025]{figdata/markets-3/06-power-markets-trading/batterycum.csv};
\legend{2024,2025}
\end{axis}
\end{tikzpicture}

\omcaption{Cumulative day-ahead arbitrage revenue of the illustrative battery (one cycle a
day, perfect foresight), per MW of power, from 1 January. Revenue accrues fastest in the
sunny, volatile months. Prices: Bundesnetzagentur | SMARD.de, CC BY 4.0.}\label{fig:m3:power-markets-trading:cum}
\end{omfigure}

\section{Tutorial: settling a day and running a battery}\label{tut:m3:power-markets-trading:battery}

\begin{tutorial}
\textbf{Goal.} Settle a producer's day and schedule a battery against day-ahead prices.
\textbf{End state:} \cref{ex:m3:power-markets-trading:wind},
\cref{fig:m3:power-markets-trading:battery,fig:m3:power-markets-trading:cum} and the annual
revenues per MW.
\begin{tutsteps}
\item \textbf{Settlement and capture.}
\omcode{code/firm/balancing/firm_balancing.py}{39}{60}{Imbalance settlement, a producer's day and the capture price.}\label{lst:m3:power-markets-trading:settle}
\item \textbf{The battery.} \Cref{met:m3:power-markets-trading:battery}.
\omcode{code/markets-3/06-power-markets-trading/python/m3_trading.py}{41}{60}{Best day-ahead schedule of a battery by dynamic programming.}\label{lst:m3:power-markets-trading:dp}
\item \textbf{Run} \texttt{m3\_trading.wind\_day()}, \texttt{capture(2025)},
\texttt{battery\_year(2025)} and \texttt{fig\_trading.py}.
\end{tutsteps}
\textbf{What to change next.} Replace perfect foresight by yesterday's prices as the forecast
and measure the revenue lost; let the battery bid 15-minute products in the auction from
October 2025.
\end{tutorial}

\section{Build: the balancing-group keeper}\label{bld:m3:power-markets-trading:balancing}

\begin{build}
\bfield{Purpose} The miniature firm runs a balancing group for its wind, solar and battery
assets: it must know, per MTU, what it has scheduled, what it metered, what it owes the grid
operator, and what its output was worth on the spot market.
\bfield{Interface} \texttt{BalancingGroup.trade(mtu, qty, price)},
\texttt{schedule\_production}, \texttt{schedule(mtu)}; \texttt{imbalance};
\texttt{settle(volume, price\_long, price\_short)}; \texttt{day\_settlement(sales, metered,
imbalance\_price)}; \texttt{capture\_price(prices, output)}.
\bfield{Rules} Positive quantities are energy into the group; a long imbalance is sold to the
operator, a short one bought from it, at a single price unless two are given; \omterm{def:m3:power-markets-trading:capture}{capture price} is
undefined without output.
\bfield{Acceptance tests} \texttt{code/firm/balancing/tests/}: a balanced producer; the signs of
imbalances and their cash; a day's settlement; a \omterm{def:m3:power-markets-trading:capture}{capture price}.
\bfield{Stretch} Quarter-hour MTUs; dual imbalance pricing; the nominations of
\cref{ch:m3:getting-access-commodities-and-power} built from this keeper.
\end{build}

\begin{omsources}
\item ENTSO-E, Single Intraday Coupling; Austrian Power Grid, ``SIDC: successful go-live of the
30-minute gate closure time in cross-zonal intraday trading''.
\item German TSOs, \emph{Description of the balancing process and the balancing markets in
Germany}; netztransparenz.de, uniform imbalance price (reBAP).
\item PJM, Financial Transmission Rights; PJM Manual 06 and the FTR forfeiture rule.
\item Bundesnetzagentur | SMARD.de, hourly prices and generation 2024--2025.
\item Pexapark, ``Unpacking H1 Deal Flow'' (2025); M.~Stover (TSSA), ``Understanding Battery Energy Storage Systems --- Current and Future'',
presentation to the ERCOT Technology and Security Committee, 9 February 2026.
\end{omsources}

\section{Exercises}

\begin{exercise}[$\star$]\label{exo:m3:power-markets-trading:1}
A solar park produces 10, 40, 40 and 10~MWh in four hours priced 80, 20, 10 and
\qty{90}{EUR/MWh}. Give its \omterm{def:m3:power-markets-trading:capture}{capture price}, the baseload price and its capture rate.
\end{exercise}

\begin{exercise}[$\star$]\label{exo:m3:power-markets-trading:2}
A BRP is 20~MWh short in a quarter-hour whose \omterm{def:m3:power-markets-trading:imbalance}{imbalance price} is \qty{300}{EUR/MWh}. What does it
pay? What would it have received had it been 20~MWh long?
\end{exercise}

\begin{exercise}[$\star$]\label{exo:m3:power-markets-trading:3}
Why does intraday trading exist at all, if the day-ahead auction already sets a price for every
hour?
\end{exercise}

\begin{exercise}[$\star\star$]\label{exo:m3:power-markets-trading:4}
An FTR of 50~MW from node $A$ to node $B$ is held for one day. The day-ahead congestion prices are
\qty{-5}{} at $A$ and \qty{+15}{EUR/MWh} at $B$ in 10 hours, and equal otherwise. What does it
pay?
\end{exercise}

\begin{exercise}[$\star\star$]\label{exo:m3:power-markets-trading:5}
In the battery of the text, what is the revenue of one day on which the two cheapest hours cost
\qty{20}{} and the two dearest pay \qty{120}{EUR/MWh}, with the dearest after the cheapest?
\end{exercise}

\begin{exercise}[$\star\star$]\label{exo:m3:power-markets-trading:6}
A buyer signs a pay-as-produced solar PPA at \qty{60}{EUR/MWh} for 100~GWh a year. Using the 2025
\omterm{def:m3:power-markets-trading:capture}{capture price}, what does the PPA cost it above the spot value of the output?
\end{exercise}

\begin{exercise}[$\star\star\star$]\label{exo:m3:power-markets-trading:7}
\emph{Coding.} With \texttt{battery\_year}, give the annual revenue per MW for 2024 and 2025 with
one cycle a day, and for 2025 with two.
\end{exercise}

\begin{exercise}[$\star\star\star$]\label{exo:m3:power-markets-trading:8}
\emph{Find the flaw.} ``Our backtest shows the battery earns EUR~75{,}000 per MW a year on
day-ahead prices; that is what we will earn.''
\end{exercise}

\section{Problem: A Battery's Year}

\begin{problem}[{Weekend problem --- 100 MW / 200 MWh in the German day-ahead market}]\label{pb:m3:power-markets-trading:1}
An investor considers a 100~MW, 200~MWh battery with a \omterm{def:m3:power-markets-trading:rte}{round-trip efficiency} of 88\%, allowed
one full cycle a day, trading only the German day-ahead auction with perfect foresight.

\medskip\noindent\textbf{Part I --- One day.}
\begin{enumerate}
\item On 18 June 2025, when does it charge and discharge, and why then?
\item How much energy does it draw from the grid, and how much does it deliver?
\item What does it earn that day?
\item Why does it charge at negative prices rather than at zero?
\item What would a second cycle have added?
\end{enumerate}

\medskip\noindent\textbf{Part II --- The year.}
\begin{enumerate}[resume]
\item Give the revenue per MW in 2024 and in 2025.
\item Why was 2025 better?
\item What does a second daily cycle add in 2025?
\item Why does revenue concentrate in some months?
\item What does the efficiency cost, roughly, on a day like 18 June?
\end{enumerate}

\medskip\noindent\textbf{Part III --- Realism.}
\begin{enumerate}[resume]
\item Why does perfect foresight overstate revenue?
\item What revenues does the calculation leave out?
\item What happens to these spreads when many batteries are built?
\item How do cycles limit the battery's life, and why cap them?
\item What would a 15-minute market change?
\end{enumerate}

\medskip\noindent\textbf{Part IV --- Judgement.}
\begin{enumerate}[resume]
\item Why do solar capture rates and battery revenues move together?
\item Who else competes for the midday cheap hours?
\item Is the day-ahead spread a durable revenue?
\item State the \emph{named result}: the battery's day-ahead arbitrage revenue per MW-year.
\item In one sentence: what does a battery sell?
\end{enumerate}
\end{problem}

\section{Interview questions}

\begin{interviewq}[$\star$ \iqroles{trader}]\label{iq:m3:power-markets-trading:1}
What is an \omterm{def:m3:power-markets-trading:imbalance}{imbalance price}, and when is being short cheap?
\end{interviewq}

\begin{interviewq}[$\star$ \iqroles{trader, researcher}]\label{iq:m3:power-markets-trading:2}
Why is the \omterm{def:m3:power-markets-trading:capture}{capture price} of solar lower than the baseload price, and why does it keep falling?
\end{interviewq}

\begin{interviewq}[$\star\star$ \iqroles{researcher}]\label{iq:m3:power-markets-trading:3}
How would you value a ten-year pay-as-produced wind PPA?
\end{interviewq}

\begin{interviewq}[$\star\star$ \iqroles{trader}]\label{iq:m3:power-markets-trading:4}
Explain a \omterm{def:m3:power-markets-trading:virtual}{virtual bid}, and how it could be used to manipulate.
\end{interviewq}

\begin{interviewq}[$\star\star$ \iqroles{risk}]\label{iq:m3:power-markets-trading:5}
A wind portfolio's forecast error is 10\% of output. What determines its expected imbalance
cost?
\end{interviewq}

\begin{interviewq}[$\star\star\star$ \iqroles{developer, researcher}]\label{iq:m3:power-markets-trading:6}
Design the optimiser that bids a battery in the day-ahead auction, the \omterm{def:m3:power-markets-trading:intraday}{intraday market} and a
reserve market at once.
\end{interviewq}
